\section{Adaptive linear flags in polynomial space}
\label{flag:section}

The strongest polynomial-space coefficient in the reviewed incoming packets
is $55.7094783189\ldots$. Its evaluator uses a fixed balanced binary tree
on sixteen groups. We improve that step by choosing a tree separately for
each exact pair-budget array. The choice is deterministic, needs only
polynomial space, and pays for every recursive child contraction.

\subsection{The finite separator problem}

Fix one exact pair-budget array. For every directed pair $e=(g,h)$,
$g\ne h$, let $\mathcal I_e$ be its finite mode-index list and put
$h_e=\log_2|\mathcal I_e|$. Arrays with an empty list contribute zero.
The exact graphical sum has the form
\[
 \sum_{(\xi_e)\in\prod_e\mathcal I_e}
       \prod_{g=1}^r Z_g((\xi_e)_{e\ni g}).
\]
Each $Z_g$ is evaluated by the prescribed child-network contraction.
For a binary separator tree $T$, let $E(T,x)$ be the union of the directed
cut-pair sets on the path from its root to leaf $x$. Successively conditioning
on the cut indices, multiplying the two conditional sums, and streaming the
outer sum is exact by finite distributivity. Its number of child-factor
evaluations is at most
\begin{equation}\label{flag:callcount}
 \sum_x\prod_{e\in E(T,x)}|\mathcal I_e|.
\end{equation}
The cut sets on a path are disjoint. Indices internal to a sibling branch
are summed independently and do not multiply the number of calls in the
chosen branch. The assignment ranks, ring accumulators, and recursion stack
have polynomial bit size.

We use the inherited weighted entropy estimate in the following notation:
for any fixed nonnegative vector $z$ on directed pairs,
\begin{equation}\label{flag:entropy}
 \sum_e z_e h_e\le
 \left(\frac{\|z\|_2}{r}+o(1)\right)K(M,u),\qquad
 K(M,u)=\frac{\pi\sqrt2}{\ln2}\frac{M}{\sqrt u}.
\end{equation}
The error is uniform on a fixed finite family of such vectors. This is the
slope-four mode entropy bound, with the sharp local precision
$(6+o(1))MD/u$ and nearly lossless asymmetric padding. It does not assume
that the child factors can be evaluated at unit cost.

\subsection{Linear flags and an intersection inequality}

Let $r=2^j$ with $j\ge2$ and identify the group labels with $V=\mathbb F_2^j$.
A complete descending linear flag
\[
 \mathcal F:\quad V=V_0>V_1>\cdots>V_j=\{0\},
 \qquad \dim V_i=j-i,
\]
defines a balanced binary tree: depth $i$ partitions $V$ into cosets of
$V_i$. Write $E(\mathcal F,x)$ for its path cut set. There are
\begin{equation}\label{flag:number}
 N_j=\prod_{i=1}^j(2^i-1)
\end{equation}
flags, since a $d$-dimensional binary vector space has $2^d-1$ hyperplanes.
Thus $N_4=315$ and $N_5=9765$. Every flag path has
\begin{equation}\label{flag:pathsize}
 v_r=\frac{r^2}{2}+\frac{r^2}{8}+\cdots+2
     =\frac{2(r^2-1)}3
\end{equation}
directed pairs.

\begin{lemma}[Common-origin intersection bound]\label{flag:intersection}
For any flags $\mathcal F,\mathcal G$ and leaves $x,y\in V$,
\[
 |E(\mathcal F,x)\cap E(\mathcal G,y)|
 \le |E(\mathcal F,0)\cap E(\mathcal G,0)|.
\]
\end{lemma}
\begin{proof}
Fix a nonzero difference $d\in V$ and consider directed pairs $(a,a+d)$.
Let $k$ be the unique index with $d\in V_{k-1}\setminus V_k$.
The pair belongs to $E(\mathcal F,x)$ exactly when
$a\in x+V_{k-1}$. If $W_{l-1}$ is the analogous subspace for
$\mathcal G$, simultaneous membership requires
\[
 a\in (x+V_{k-1})\cap(y+W_{l-1}).
\]
This affine intersection is empty or a coset of
$V_{k-1}\cap W_{l-1}$. Its cardinality is therefore at most the
cardinality of that homogeneous intersection, attained when $x=y=0$.
Sum this inequality over $d\ne0$.
\end{proof}

\begin{lemma}[Uniform flag average]\label{flag:average}
The vector
$a_e=N_j^{-1}\sum_{\mathcal F}\mathbf1_{e\in E(\mathcal F,0)}$
satisfies
\[
 a_e=\begin{cases}
 1,&\text{one endpoint of $e$ is zero},\\
 2/3,&\text{both endpoints are nonzero and distinct}.
 \end{cases}
\]
Consequently
\begin{equation}\label{flag:norm}
 \|a\|_2^2
 =2(r-1)+\frac49(r-1)(r-2)
 =\frac{(r-1)(4r+10)}9.
\end{equation}
\end{lemma}
\begin{proof}
Every path to zero separates every directed pair incident with zero.
The group $\mathrm{GL}(j,2)$ permutes the flags and acts transitively on
ordered pairs of distinct nonzero vectors. Hence the second average value
is constant. Subtracting the $2(r-1)$ incident pairs from
\eqref{flag:pathsize} and dividing by $(r-1)(r-2)$ gives $2/3$.
Counting the two types of pairs proves \eqref{flag:norm}.
\end{proof}

\begin{theorem}[Adaptive flag separator]\label{flag:local}
For each exact pair-budget array, choose a flag minimizing
$\max_x\sum_{e\in E(\mathcal F,x)}h_e$. Streamed evaluation with that
flag uses polynomial space and obeys the paid recurrence
\begin{equation}\label{flag:recurrence}
 T_d\le\operatorname{poly}(s)
 2^{(b_{r_d}+o(1))K(M_d,u_d)}T_{d+1},\qquad
 b_r=\frac1r\sqrt{\frac{(r-1)(4r+10)}9}.
\end{equation}
A single fixed polynomial in the input-description length and $s$ pays
for input access and arithmetic, as in the common contraction interface.
\end{theorem}
\begin{proof}
For each flag select a maximizing leaf $x_{\mathcal F}$ and put
\[
 z=\frac1{N_j}\sum_{\mathcal F}
        \mathbf1_{E(\mathcal F,x_{\mathcal F})}.
\]
The minimum maximum path cost is at most its average over flags, namely
$\sum_e z_eh_e$. Expanding the squared norm and applying
Lemma~\ref{flag:intersection} to every pair of flags yields
\[
 \|z\|_2^2
 =\frac1{N_j^2}\sum_{\mathcal F,\mathcal G}
 |E(\mathcal F,x_{\mathcal F})\cap E(\mathcal G,x_{\mathcal G})|
 \le\frac1{N_j^2}\sum_{\mathcal F,\mathcal G}
 |E(\mathcal F,0)\cap E(\mathcal G,0)|
 =\|a\|_2^2.
\]
Use \eqref{flag:entropy} and Lemma~\ref{flag:average}. The vectors $z$
belong to a finite family determined only by the fixed group alphabet,
so their data-dependent choice preserves the uniform error bound.
Equation~\eqref{flag:callcount} contributes only a constant factor $r$
beyond the maximum path count and pays for the real child contraction.

The flag selection requires no logarithmic approximation. For each path,
form the exact integer product of its list sizes; for each flag take the
largest such product and select a flag with the smallest maximum. There
are only $N_jr$ products, each of polynomial bit length. Exact list indexing,
flag selection, and both separator and network recursion use polynomial
space when streamed depth first.
\end{proof}

At $r=16$, this replaces $\sqrt{170}/16$ by
$\sqrt{370/3}/16$. Keeping the incoming packet's schedule already lowers
its coefficient to $47.4509449603\ldots$. We give a slightly stronger
complete rational schedule next.

\subsection{The integer schedule and its complete bound}

\begin{theorem}[Polynomial-space certificate]\label{flag:global}
Under the common normalized contraction interface, adaptive linear flags
give deterministic polynomial-space bit time
\[
 (s+\mathcal L)^{O(1)}
 2^{(C_{\mathrm{flag}}+o(1))s/\sqrt{\log_2s}},
\]
where
\begin{equation}\label{flag:constant}
 C_{\mathrm{flag}}
 =\frac{\pi\sqrt2}{\ln2}
 \left(\frac{9\sqrt{80}}8 b_{16}
       +\frac{\sqrt{120}}{64}b_8
       +\frac{\sqrt{60}\,b_4}
        {256\sqrt{1-q}\,(1-1/(2\sqrt q))}\right)
 =46.18484558673\ldots,
\end{equation}
where $q=63/100$ and $b_r=r^{-1}\sqrt{(r-1)(4r+10)/9}$.
In particular, $46.19$ is a valid coefficient for all sufficiently large
$s$. The result retains the common exact contraction and SAT-reduction
dependencies; it does not assert a smaller logarithmic order.
\end{theorem}
\begin{proof}
Let $\lambda=\lceil\log_2s\rceil$,
$g=\lceil\sqrt\lambda\rceil$, $\Lambda=\lambda-g$, and
$b=\lceil6\log_2\lambda\rceil$. Choose
\begin{align*}
 r_0=r_1&=16,&u_0=u_1&=\left\lfloor\frac{\Lambda}{80}\right\rfloor,\\
 r_2&=8,&u_2&=\left\lfloor\frac{\Lambda}{120}\right\rfloor,\\
 r_d&=4\quad(d\ge3),&u_d&=
 \left\lfloor\frac{37\Lambda}{6000}
          \left(\frac{63}{100}\right)^{d-3}\right\rfloor\quad(d\ge3),\\
 a_d&=\sum_{q=0}^d u_q+(d+1)b,&
 H&=\left\lceil\frac{s}{\lambda}\right\rceil.
\end{align*}
The cumulative retained ranges through depths one and two are at most
$\Lambda/40$ and $\Lambda/30$, respectively. The final tail has sum at
most $\Lambda/60$, so the total is at most $\Lambda/20$.
Use direct contraction
at $M\le H$ and at bounded sizes for which the parameter or admissibility
checks fail. The following argument proves that all checks eventually pass.

First consider hypothetical active depths
$d\le\lceil\log_{20/11}\lambda\rceil+O(1)$ and $M>H$. Their minimum taper is
$\lambda^{1+\ln(63/100)/\ln(20/11)+o(1)}$, which diverges because
$63/100>11/20$. Also $a_d\le\Lambda/20+O((\log\lambda)^2)$.
At the root, $Q=16^{a_0}$ and the packing term $Q^2\Delta$ has logarithm at most
\[
 \lambda-\frac{g}{10}+O(\log\lambda),
 \qquad \Delta=\lceil s^{9/10}\rceil.
\]
At depth one $Q=16^{a_1}$, at depth two $Q=8^{a_2}$, and thereafter
$Q=4^{a_d}$. The retained side has degree at most $\Delta$, and the fresh
side has degree at most $2^{a_{d-1}+1}$. The cumulative range bounds give,
respectively, $4/40=3/30=2/20=1/10$ as the leading exponent of $Q$.
Thus the logarithm of $Q\Delta$ is always at most
$\lambda-g/10+O((\log\lambda)^2)$.
The fresh-side packing and packet-rounding exponents are at most
$(9/40)\lambda+o(\lambda)$ at depth one,
$(7/30)\lambda+o(\lambda)$ at depth two, and
$\lambda/4+o(\lambda)$ thereafter. Thus all three asymmetric padding gates hold
uniformly: the negative $\sqrt\lambda$ margin absorbs the buffers,
polylogarithmic factors, and $s/M\le\lambda$. Newly created degrees stay
below $s^{1/20+o(1)}<\Delta$.

Padding is consequently $(1+o(1))M$. The first two child ratios are
$(1+o(1))/8$, the next is $(1+o(1))/4$, and the subsequent ratio is
$(1+o(1))/2$, eventually at most $11/20$ in the tail.
This validates the hypothetical depth bound and closes the
induction. More sharply,
\[
 M_0\le s,\quad M_1\le(1+o(1))\frac{s}{8},\quad
 M_2\le(1+o(1))^2\frac{s}{64},\quad
 M_d\le(1+o(1))^d\frac{s}{256\,2^{d-3}}\quad(d\ge3).
\]
The padding error is $O(1/\lambda)$ plus terms exponentially small in
$\sqrt\lambda$ or a positive power of $s$, so its product over these
$O(\log\lambda)$ levels is $1+o(1)$.

The buffer ensures $D\ge2^{a_{d-1}+b+1}$. The inherited precision bound
$p=O(s\lambda^2 2^{a_{d-1}})$ then gives
$p/(MD/u)=O(\lambda^4 2^{-b})=o(1)$; the remaining recovery term is also
lower order. This validates the uniform local precision and entropy
estimates used in \eqref{flag:recurrence}. Floors are uniformly relative
$o(1)$ because the minimum taper diverges.

Summing the paid exponent in \eqref{flag:recurrence} gives, on the scale
$s/\sqrt\lambda$, the coefficient
\[
 \frac{\pi\sqrt2}{\ln2}
 \left(\frac{9\sqrt{80}}8 b_{16}
       +\frac{\sqrt{120}}{64}b_8
       +\frac{\sqrt{60}\,b_4}
        {256\sqrt{1-q}\,(1-1/(2\sqrt q))}\right)+o(1).
\]
Indeed the first two levels contribute
$\sqrt{80}b_{16}(1+1/8)$, and depth two contributes
$\sqrt{120}b_8/64$. In the tail, $M_d/\sqrt{u_d}$ has ratio
$1/(2\sqrt q)<1$ and first coefficient
$\sqrt{60}/(256\sqrt{1-q})$. Summing its geometric series yields
\eqref{flag:constant}.

Polynomial work over $O(\log\lambda)$ levels adds
$O(\lambda\log\lambda)$ to the time exponent. Leaves cost
$2^{O(s/\lambda)}$; original entries cost $2^{O(s^{9/10})}$; and fresh
entry routines cost at most
$\exp(O(2^{a_d}\log s))=\exp(s^{1/20+o(1)})$. These are all lower order,
including after the ancestor evaluation counts have been paid. Creation-time
tensor descriptions remain immutable under the imported conventions, so no
unaccounted recursive entry expansion occurs. All ring dimensions and
precisions are polynomial, and the exact flag choices and contractions are
streamed. This proves the bit-time and polynomial-space statements.
\end{proof}

The new coefficient is approximately $17.10\%$ below the reviewed
$55.7094783189\ldots$ polynomial-space coefficient. This compares normalized
upper-bound functions. It is not a runtime measurement and does not identify
the wire parameter $s$ with the number of SAT variables.

\subsection{A scoped optimality statement}

For a fixed $r=2^j$, the local constant $b_r$ is optimal for a choice among
balanced binary separator trees when only the Euclidean pair-entropy
constraint is used. To see this, take $h$ proportional to the average
vector $a$ in Lemma~\ref{flag:average}, with distinguished vertex zero.
Every balanced tree's path to zero contains all $2(r-1)$ pairs incident
with zero and has the same total size $v_r$. Its $a$-weight is therefore
$2(r-1)+(2/3)(v_r-2(r-1))=\|a\|_2^2$.
After normalization this matches the proved coefficient. This is a barrier
only for this restricted separator model and entropy relaxation. It is not
a lower bound for unbalanced separator trees, arbitrary contraction methods,
or SAT algorithms.
